% ---------------------------------------------------------------------------
% Appendix: source listings for the System Serial Bus Design report.
%
% 1. Add these two lines to the PREAMBLE of main.tex, if not already there:
%
%        \usepackage{listings}
%        \usepackage{xcolor}
%
% 2. Add this line at the very end of main.tex, after Section 13 (Conclusion)
%    and before \end{document}:
%
%        % ---------------------------------------------------------------------------
% Appendix: source listings for the System Serial Bus Design report.
%
% 1. Add these two lines to the PREAMBLE of main.tex, if not already there:
%
%        \usepackage{listings}
%        \usepackage{xcolor}
%
% 2. Add this line at the very end of main.tex, after Section 13 (Conclusion)
%    and before \end{document}:
%
%        % ---------------------------------------------------------------------------
% Appendix: source listings for the System Serial Bus Design report.
%
% 1. Add these two lines to the PREAMBLE of main.tex, if not already there:
%
%        \usepackage{listings}
%        \usepackage{xcolor}
%
% 2. Add this line at the very end of main.tex, after Section 13 (Conclusion)
%    and before \end{document}:
%
%        % ---------------------------------------------------------------------------
% Appendix: source listings for the System Serial Bus Design report.
%
% 1. Add these two lines to the PREAMBLE of main.tex, if not already there:
%
%        \usepackage{listings}
%        \usepackage{xcolor}
%
% 2. Add this line at the very end of main.tex, after Section 13 (Conclusion)
%    and before \end{document}:
%
%        \input{appendix}
%
% 3. Upload the whole appendix/ folder next to main.tex.
%
% The page border is drawn by main.tex, so these pages inherit it automatically.
% Every listing has had its comments removed so it prints compactly; the code
% is otherwise unchanged and still compiles.
% ---------------------------------------------------------------------------

\definecolor{kwcol}{RGB}{0,0,140}
\definecolor{cmtcol}{RGB}{90,90,90}
\definecolor{strcol}{RGB}{140,0,0}
\definecolor{numcol}{RGB}{140,140,140}

\lstdefinelanguage{SystemVerilog}[]{Verilog}{%
  morekeywords={logic,bit,byte,int,longint,shortint,void,always_ff,always_comb,%
    always_latch,typedef,enum,struct,union,packed,unsigned,signed,localparam,%
    interface,endinterface,modport,package,endpackage,import,export,program,%
    endprogram,property,endproperty,sequence,endsequence,assert,assume,cover,%
    covergroup,endgroup,coverpoint,constraint,rand,randc,unique,priority,%
    automatic,static,ref,const,break,continue,return,do,final,foreach,inside,%
    iff,with,timeunit,timeprecision,generate,endgenerate,genvar,string,%
    this,super,extends,virtual,pure,class,endclass,new,null},%
  sensitive=true,%
}

\lstdefinestyle{codeprint}{%
  basicstyle=\ttfamily\scriptsize,
  keywordstyle=\bfseries\color{kwcol},
  commentstyle=\itshape\color{cmtcol},
  stringstyle=\color{strcol},
  numbers=left,
  numberstyle=\tiny\color{numcol},
  numbersep=7pt,
  stepnumber=1,
  breaklines=true,
  breakatwhitespace=false,
  postbreak=\mbox{\textcolor{numcol}{$\hookrightarrow$}\space},
  showstringspaces=false,
  tabsize=4,
  columns=fullflexible,
  keepspaces=true,
  upquote=true,
  alsoletter=_,
  frame=none,
  xleftmargin=2.4em,
  aboveskip=0.4em,
  belowskip=1.0em,
}
\lstset{style=codeprint}

% #1 real filename, #2 filename with escaped underscores, #3 caption, #4 language
\newcommand{\srcentry}[4]{%
  \subsection{\texttt{#2}}%
  \noindent\small #3\par\medskip\nopagebreak
  \lstinputlisting[language=#4]{appendix/#1}%
}

\clearpage
\appendix
\phantomsection
\section*{Appendix: Source Code}
\addcontentsline{toc}{section}{Appendix: Source Code}

\noindent
The listings that follow are the complete source of the design described above,
presented in the order the report introduces it. Comments have been removed so
that the code prints compactly; the sources are otherwise unchanged and compile
as shown.

\medskip

\section{Top Level and System Assembly}
\srcentry{Sytembus_top.sv}{Sytembus\_top.sv}{JTAG wrapper: In-System Sources and Probes driving both master device ports and the remote link pins (Section 11).}{SystemVerilog}
\srcentry{Sytembus.sv}{Sytembus.sv}{System assembly: two master ports, three slaves, the bridge node and the interconnect (Figure 1).}{SystemVerilog}

\clearpage
\section{Bus Interconnect}
\srcentry{bus_m3_s4.sv}{bus\_m3\_s4.sv}{Interconnect of Section 3: arbiter, decoder, multiplexers, acknowledge gating and second-split blocking.}{SystemVerilog}
\srcentry{muxes.sv}{muxes.sv}{The 3-to-1 and 4-to-1 single-bit multiplexers of Table 12.}{SystemVerilog}

\clearpage
\section{Bus Bridge}
\srcentry{bus_bridge_node.sv}{bus\_bridge\_node.sv}{Bridge slave face, master face, client and server state machines, and response timeout (Sections 4.2 to 4.8).}{SystemVerilog}
\srcentry{uart.sv}{uart.sv}{8N1 transmitter and receiver at 115200 baud, 434 clocks per bit (Section 4.4).}{SystemVerilog}

\clearpage
\section{Arbiter}
\srcentry{arbiter3.sv}{arbiter3.sv}{Fixed-priority arbitration, split ownership and ready gating (Sections 5.3 to 5.5).}{SystemVerilog}

\clearpage
\section{Master Port and Slave Port}
\srcentry{master_port.sv}{master\_port.sv}{Parallel to serial: address and data shifted out one bit per clock (Section 6.1).}{SystemVerilog}
\srcentry{slave_port.sv}{slave\_port.sv}{Serial to parallel: offset and data collected, read data shifted back (Section 6.1).}{SystemVerilog}
\srcentry{slave.sv}{slave.sv}{Slave wrapper pairing a slave port with its block RAM.}{SystemVerilog}
\srcentry{slave_memory_bram.sv}{slave\_memory\_bram.sv}{Inferred single-port block RAM used by all three slaves (Section 11).}{SystemVerilog}

\clearpage
\section{Address Decoder}
\srcentry{addr_decoder4.sv}{addr\_decoder4.sv}{Device-ID decode, valid routing, acknowledge, and per-target split detection (Sections 8.3 to 8.5).}{SystemVerilog}

\clearpage
\section{Timing Constraints}
\srcentry{Sytembus.sdc}{Sytembus.sdc}{Clock definition and the reset false path quoted in Section 10.}{tcl}

\clearpage
\section{Simulation and Verification}
\srcentry{final_tb.sv}{final\_tb.sv}{System testbench of Sections 12.2 to 12.8: reset, single-master transfers, arbitration, split, and decode error.}{SystemVerilog}
\srcentry{arbiter_tb.sv}{arbiter\_tb.sv}{Unit-level arbiter checks (Section 12.1).}{SystemVerilog}
\srcentry{addr_decoder_tb.sv}{addr\_decoder\_tb.sv}{Unit-level address decoder checks (Section 12.1).}{SystemVerilog}
\srcentry{master2_slave3_tb.sv}{master2\_slave3\_tb.sv}{Two-master and split-capable slave integration test (Sections 12.5 and 12.6).}{SystemVerilog}
\srcentry{remote_link_tb.sv}{remote\_link\_tb.sv}{Bridge round trip over the UART link (Section 4).}{SystemVerilog}
\srcentry{frames_tb.sv}{frames\_tb.sv}{Wire-format checks against the frames of Tables 14 to 16.}{SystemVerilog}
\srcentry{figures_tb.sv}{figures\_tb.sv}{Generates the waveform time windows recorded in figures.txt (Section 12.2).}{SystemVerilog}

\clearpage
\section{Build and Hardware Test Scripts}
\srcentry{scr.tcl}{scr.tcl}{ModelSim driver: compiles the design and runs each testbench (Section 12).}{tcl}
\srcentry{issp_bus_test.tcl}{issp\_bus\_test.tcl}{Scripted hardware verification over JTAG (Section 11).}{tcl}
\srcentry{issp_console.tcl}{issp\_console.tcl}{Interactive memory console over In-System Sources and Probes.}{tcl}

%
% 3. Upload the whole appendix/ folder next to main.tex.
%
% The page border is drawn by main.tex, so these pages inherit it automatically.
% Every listing has had its comments removed so it prints compactly; the code
% is otherwise unchanged and still compiles.
% ---------------------------------------------------------------------------

\definecolor{kwcol}{RGB}{0,0,140}
\definecolor{cmtcol}{RGB}{90,90,90}
\definecolor{strcol}{RGB}{140,0,0}
\definecolor{numcol}{RGB}{140,140,140}

\lstdefinelanguage{SystemVerilog}[]{Verilog}{%
  morekeywords={logic,bit,byte,int,longint,shortint,void,always_ff,always_comb,%
    always_latch,typedef,enum,struct,union,packed,unsigned,signed,localparam,%
    interface,endinterface,modport,package,endpackage,import,export,program,%
    endprogram,property,endproperty,sequence,endsequence,assert,assume,cover,%
    covergroup,endgroup,coverpoint,constraint,rand,randc,unique,priority,%
    automatic,static,ref,const,break,continue,return,do,final,foreach,inside,%
    iff,with,timeunit,timeprecision,generate,endgenerate,genvar,string,%
    this,super,extends,virtual,pure,class,endclass,new,null},%
  sensitive=true,%
}

\lstdefinestyle{codeprint}{%
  basicstyle=\ttfamily\scriptsize,
  keywordstyle=\bfseries\color{kwcol},
  commentstyle=\itshape\color{cmtcol},
  stringstyle=\color{strcol},
  numbers=left,
  numberstyle=\tiny\color{numcol},
  numbersep=7pt,
  stepnumber=1,
  breaklines=true,
  breakatwhitespace=false,
  postbreak=\mbox{\textcolor{numcol}{$\hookrightarrow$}\space},
  showstringspaces=false,
  tabsize=4,
  columns=fullflexible,
  keepspaces=true,
  upquote=true,
  alsoletter=_,
  frame=none,
  xleftmargin=2.4em,
  aboveskip=0.4em,
  belowskip=1.0em,
}
\lstset{style=codeprint}

% #1 real filename, #2 filename with escaped underscores, #3 caption, #4 language
\newcommand{\srcentry}[4]{%
  \subsection{\texttt{#2}}%
  \noindent\small #3\par\medskip\nopagebreak
  \lstinputlisting[language=#4]{appendix/#1}%
}

\clearpage
\appendix
\phantomsection
\section*{Appendix: Source Code}
\addcontentsline{toc}{section}{Appendix: Source Code}

\noindent
The listings that follow are the complete source of the design described above,
presented in the order the report introduces it. Comments have been removed so
that the code prints compactly; the sources are otherwise unchanged and compile
as shown.

\medskip

\section{Top Level and System Assembly}
\srcentry{Sytembus_top.sv}{Sytembus\_top.sv}{JTAG wrapper: In-System Sources and Probes driving both master device ports and the remote link pins (Section 11).}{SystemVerilog}
\srcentry{Sytembus.sv}{Sytembus.sv}{System assembly: two master ports, three slaves, the bridge node and the interconnect (Figure 1).}{SystemVerilog}

\clearpage
\section{Bus Interconnect}
\srcentry{bus_m3_s4.sv}{bus\_m3\_s4.sv}{Interconnect of Section 3: arbiter, decoder, multiplexers, acknowledge gating and second-split blocking.}{SystemVerilog}
\srcentry{muxes.sv}{muxes.sv}{The 3-to-1 and 4-to-1 single-bit multiplexers of Table 12.}{SystemVerilog}

\clearpage
\section{Bus Bridge}
\srcentry{bus_bridge_node.sv}{bus\_bridge\_node.sv}{Bridge slave face, master face, client and server state machines, and response timeout (Sections 4.2 to 4.8).}{SystemVerilog}
\srcentry{uart.sv}{uart.sv}{8N1 transmitter and receiver at 115200 baud, 434 clocks per bit (Section 4.4).}{SystemVerilog}

\clearpage
\section{Arbiter}
\srcentry{arbiter3.sv}{arbiter3.sv}{Fixed-priority arbitration, split ownership and ready gating (Sections 5.3 to 5.5).}{SystemVerilog}

\clearpage
\section{Master Port and Slave Port}
\srcentry{master_port.sv}{master\_port.sv}{Parallel to serial: address and data shifted out one bit per clock (Section 6.1).}{SystemVerilog}
\srcentry{slave_port.sv}{slave\_port.sv}{Serial to parallel: offset and data collected, read data shifted back (Section 6.1).}{SystemVerilog}
\srcentry{slave.sv}{slave.sv}{Slave wrapper pairing a slave port with its block RAM.}{SystemVerilog}
\srcentry{slave_memory_bram.sv}{slave\_memory\_bram.sv}{Inferred single-port block RAM used by all three slaves (Section 11).}{SystemVerilog}

\clearpage
\section{Address Decoder}
\srcentry{addr_decoder4.sv}{addr\_decoder4.sv}{Device-ID decode, valid routing, acknowledge, and per-target split detection (Sections 8.3 to 8.5).}{SystemVerilog}

\clearpage
\section{Timing Constraints}
\srcentry{Sytembus.sdc}{Sytembus.sdc}{Clock definition and the reset false path quoted in Section 10.}{tcl}

\clearpage
\section{Simulation and Verification}
\srcentry{final_tb.sv}{final\_tb.sv}{System testbench of Sections 12.2 to 12.8: reset, single-master transfers, arbitration, split, and decode error.}{SystemVerilog}
\srcentry{arbiter_tb.sv}{arbiter\_tb.sv}{Unit-level arbiter checks (Section 12.1).}{SystemVerilog}
\srcentry{addr_decoder_tb.sv}{addr\_decoder\_tb.sv}{Unit-level address decoder checks (Section 12.1).}{SystemVerilog}
\srcentry{master2_slave3_tb.sv}{master2\_slave3\_tb.sv}{Two-master and split-capable slave integration test (Sections 12.5 and 12.6).}{SystemVerilog}
\srcentry{remote_link_tb.sv}{remote\_link\_tb.sv}{Bridge round trip over the UART link (Section 4).}{SystemVerilog}
\srcentry{frames_tb.sv}{frames\_tb.sv}{Wire-format checks against the frames of Tables 14 to 16.}{SystemVerilog}
\srcentry{figures_tb.sv}{figures\_tb.sv}{Generates the waveform time windows recorded in figures.txt (Section 12.2).}{SystemVerilog}

\clearpage
\section{Build and Hardware Test Scripts}
\srcentry{scr.tcl}{scr.tcl}{ModelSim driver: compiles the design and runs each testbench (Section 12).}{tcl}
\srcentry{issp_bus_test.tcl}{issp\_bus\_test.tcl}{Scripted hardware verification over JTAG (Section 11).}{tcl}
\srcentry{issp_console.tcl}{issp\_console.tcl}{Interactive memory console over In-System Sources and Probes.}{tcl}

%
% 3. Upload the whole appendix/ folder next to main.tex.
%
% The page border is drawn by main.tex, so these pages inherit it automatically.
% Every listing has had its comments removed so it prints compactly; the code
% is otherwise unchanged and still compiles.
% ---------------------------------------------------------------------------

\definecolor{kwcol}{RGB}{0,0,140}
\definecolor{cmtcol}{RGB}{90,90,90}
\definecolor{strcol}{RGB}{140,0,0}
\definecolor{numcol}{RGB}{140,140,140}

\lstdefinelanguage{SystemVerilog}[]{Verilog}{%
  morekeywords={logic,bit,byte,int,longint,shortint,void,always_ff,always_comb,%
    always_latch,typedef,enum,struct,union,packed,unsigned,signed,localparam,%
    interface,endinterface,modport,package,endpackage,import,export,program,%
    endprogram,property,endproperty,sequence,endsequence,assert,assume,cover,%
    covergroup,endgroup,coverpoint,constraint,rand,randc,unique,priority,%
    automatic,static,ref,const,break,continue,return,do,final,foreach,inside,%
    iff,with,timeunit,timeprecision,generate,endgenerate,genvar,string,%
    this,super,extends,virtual,pure,class,endclass,new,null},%
  sensitive=true,%
}

\lstdefinestyle{codeprint}{%
  basicstyle=\ttfamily\scriptsize,
  keywordstyle=\bfseries\color{kwcol},
  commentstyle=\itshape\color{cmtcol},
  stringstyle=\color{strcol},
  numbers=left,
  numberstyle=\tiny\color{numcol},
  numbersep=7pt,
  stepnumber=1,
  breaklines=true,
  breakatwhitespace=false,
  postbreak=\mbox{\textcolor{numcol}{$\hookrightarrow$}\space},
  showstringspaces=false,
  tabsize=4,
  columns=fullflexible,
  keepspaces=true,
  upquote=true,
  alsoletter=_,
  frame=none,
  xleftmargin=2.4em,
  aboveskip=0.4em,
  belowskip=1.0em,
}
\lstset{style=codeprint}

% #1 real filename, #2 filename with escaped underscores, #3 caption, #4 language
\newcommand{\srcentry}[4]{%
  \subsection{\texttt{#2}}%
  \noindent\small #3\par\medskip\nopagebreak
  \lstinputlisting[language=#4]{appendix/#1}%
}

\clearpage
\appendix
\phantomsection
\section*{Appendix: Source Code}
\addcontentsline{toc}{section}{Appendix: Source Code}

\noindent
The listings that follow are the complete source of the design described above,
presented in the order the report introduces it. Comments have been removed so
that the code prints compactly; the sources are otherwise unchanged and compile
as shown.

\medskip

\section{Top Level and System Assembly}
\srcentry{Sytembus_top.sv}{Sytembus\_top.sv}{JTAG wrapper: In-System Sources and Probes driving both master device ports and the remote link pins (Section 11).}{SystemVerilog}
\srcentry{Sytembus.sv}{Sytembus.sv}{System assembly: two master ports, three slaves, the bridge node and the interconnect (Figure 1).}{SystemVerilog}

\clearpage
\section{Bus Interconnect}
\srcentry{bus_m3_s4.sv}{bus\_m3\_s4.sv}{Interconnect of Section 3: arbiter, decoder, multiplexers, acknowledge gating and second-split blocking.}{SystemVerilog}
\srcentry{muxes.sv}{muxes.sv}{The 3-to-1 and 4-to-1 single-bit multiplexers of Table 12.}{SystemVerilog}

\clearpage
\section{Bus Bridge}
\srcentry{bus_bridge_node.sv}{bus\_bridge\_node.sv}{Bridge slave face, master face, client and server state machines, and response timeout (Sections 4.2 to 4.8).}{SystemVerilog}
\srcentry{uart.sv}{uart.sv}{8N1 transmitter and receiver at 115200 baud, 434 clocks per bit (Section 4.4).}{SystemVerilog}

\clearpage
\section{Arbiter}
\srcentry{arbiter3.sv}{arbiter3.sv}{Fixed-priority arbitration, split ownership and ready gating (Sections 5.3 to 5.5).}{SystemVerilog}

\clearpage
\section{Master Port and Slave Port}
\srcentry{master_port.sv}{master\_port.sv}{Parallel to serial: address and data shifted out one bit per clock (Section 6.1).}{SystemVerilog}
\srcentry{slave_port.sv}{slave\_port.sv}{Serial to parallel: offset and data collected, read data shifted back (Section 6.1).}{SystemVerilog}
\srcentry{slave.sv}{slave.sv}{Slave wrapper pairing a slave port with its block RAM.}{SystemVerilog}
\srcentry{slave_memory_bram.sv}{slave\_memory\_bram.sv}{Inferred single-port block RAM used by all three slaves (Section 11).}{SystemVerilog}

\clearpage
\section{Address Decoder}
\srcentry{addr_decoder4.sv}{addr\_decoder4.sv}{Device-ID decode, valid routing, acknowledge, and per-target split detection (Sections 8.3 to 8.5).}{SystemVerilog}

\clearpage
\section{Timing Constraints}
\srcentry{Sytembus.sdc}{Sytembus.sdc}{Clock definition and the reset false path quoted in Section 10.}{tcl}

\clearpage
\section{Simulation and Verification}
\srcentry{final_tb.sv}{final\_tb.sv}{System testbench of Sections 12.2 to 12.8: reset, single-master transfers, arbitration, split, and decode error.}{SystemVerilog}
\srcentry{arbiter_tb.sv}{arbiter\_tb.sv}{Unit-level arbiter checks (Section 12.1).}{SystemVerilog}
\srcentry{addr_decoder_tb.sv}{addr\_decoder\_tb.sv}{Unit-level address decoder checks (Section 12.1).}{SystemVerilog}
\srcentry{master2_slave3_tb.sv}{master2\_slave3\_tb.sv}{Two-master and split-capable slave integration test (Sections 12.5 and 12.6).}{SystemVerilog}
\srcentry{remote_link_tb.sv}{remote\_link\_tb.sv}{Bridge round trip over the UART link (Section 4).}{SystemVerilog}
\srcentry{frames_tb.sv}{frames\_tb.sv}{Wire-format checks against the frames of Tables 14 to 16.}{SystemVerilog}
\srcentry{figures_tb.sv}{figures\_tb.sv}{Generates the waveform time windows recorded in figures.txt (Section 12.2).}{SystemVerilog}

\clearpage
\section{Build and Hardware Test Scripts}
\srcentry{scr.tcl}{scr.tcl}{ModelSim driver: compiles the design and runs each testbench (Section 12).}{tcl}
\srcentry{issp_bus_test.tcl}{issp\_bus\_test.tcl}{Scripted hardware verification over JTAG (Section 11).}{tcl}
\srcentry{issp_console.tcl}{issp\_console.tcl}{Interactive memory console over In-System Sources and Probes.}{tcl}

%
% 3. Upload the whole appendix/ folder next to main.tex.
%
% The page border is drawn by main.tex, so these pages inherit it automatically.
% Every listing has had its comments removed so it prints compactly; the code
% is otherwise unchanged and still compiles.
% ---------------------------------------------------------------------------

\definecolor{kwcol}{RGB}{0,0,140}
\definecolor{cmtcol}{RGB}{90,90,90}
\definecolor{strcol}{RGB}{140,0,0}
\definecolor{numcol}{RGB}{140,140,140}

\lstdefinelanguage{SystemVerilog}[]{Verilog}{%
  morekeywords={logic,bit,byte,int,longint,shortint,void,always_ff,always_comb,%
    always_latch,typedef,enum,struct,union,packed,unsigned,signed,localparam,%
    interface,endinterface,modport,package,endpackage,import,export,program,%
    endprogram,property,endproperty,sequence,endsequence,assert,assume,cover,%
    covergroup,endgroup,coverpoint,constraint,rand,randc,unique,priority,%
    automatic,static,ref,const,break,continue,return,do,final,foreach,inside,%
    iff,with,timeunit,timeprecision,generate,endgenerate,genvar,string,%
    this,super,extends,virtual,pure,class,endclass,new,null},%
  sensitive=true,%
}

\lstdefinestyle{codeprint}{%
  basicstyle=\ttfamily\scriptsize,
  keywordstyle=\bfseries\color{kwcol},
  commentstyle=\itshape\color{cmtcol},
  stringstyle=\color{strcol},
  numbers=left,
  numberstyle=\tiny\color{numcol},
  numbersep=7pt,
  stepnumber=1,
  breaklines=true,
  breakatwhitespace=false,
  postbreak=\mbox{\textcolor{numcol}{$\hookrightarrow$}\space},
  showstringspaces=false,
  tabsize=4,
  columns=fullflexible,
  keepspaces=true,
  upquote=true,
  alsoletter=_,
  frame=none,
  xleftmargin=2.4em,
  aboveskip=0.4em,
  belowskip=1.0em,
}
\lstset{style=codeprint}

% #1 real filename, #2 filename with escaped underscores, #3 caption, #4 language
\newcommand{\srcentry}[4]{%
  \subsection{\texttt{#2}}%
  \noindent\small #3\par\medskip\nopagebreak
  \lstinputlisting[language=#4]{appendix/#1}%
}

\clearpage
\appendix
\phantomsection
\section*{Appendix: Source Code}
\addcontentsline{toc}{section}{Appendix: Source Code}

\noindent
The listings that follow are the complete source of the design described above,
presented in the order the report introduces it. Comments have been removed so
that the code prints compactly; the sources are otherwise unchanged and compile
as shown.

\medskip

\section{Top Level and System Assembly}
\srcentry{Sytembus_top.sv}{Sytembus\_top.sv}{JTAG wrapper: In-System Sources and Probes driving both master device ports and the remote link pins (Section 11).}{SystemVerilog}
\srcentry{Sytembus.sv}{Sytembus.sv}{System assembly: two master ports, three slaves, the bridge node and the interconnect (Figure 1).}{SystemVerilog}

\clearpage
\section{Bus Interconnect}
\srcentry{bus_m3_s4.sv}{bus\_m3\_s4.sv}{Interconnect of Section 3: arbiter, decoder, multiplexers, acknowledge gating and second-split blocking.}{SystemVerilog}
\srcentry{muxes.sv}{muxes.sv}{The 3-to-1 and 4-to-1 single-bit multiplexers of Table 12.}{SystemVerilog}

\clearpage
\section{Bus Bridge}
\srcentry{bus_bridge_node.sv}{bus\_bridge\_node.sv}{Bridge slave face, master face, client and server state machines, and response timeout (Sections 4.2 to 4.8).}{SystemVerilog}
\srcentry{uart.sv}{uart.sv}{8N1 transmitter and receiver at 115200 baud, 434 clocks per bit (Section 4.4).}{SystemVerilog}

\clearpage
\section{Arbiter}
\srcentry{arbiter3.sv}{arbiter3.sv}{Fixed-priority arbitration, split ownership and ready gating (Sections 5.3 to 5.5).}{SystemVerilog}

\clearpage
\section{Master Port and Slave Port}
\srcentry{master_port.sv}{master\_port.sv}{Parallel to serial: address and data shifted out one bit per clock (Section 6.1).}{SystemVerilog}
\srcentry{slave_port.sv}{slave\_port.sv}{Serial to parallel: offset and data collected, read data shifted back (Section 6.1).}{SystemVerilog}
\srcentry{slave.sv}{slave.sv}{Slave wrapper pairing a slave port with its block RAM.}{SystemVerilog}
\srcentry{slave_memory_bram.sv}{slave\_memory\_bram.sv}{Inferred single-port block RAM used by all three slaves (Section 11).}{SystemVerilog}

\clearpage
\section{Address Decoder}
\srcentry{addr_decoder4.sv}{addr\_decoder4.sv}{Device-ID decode, valid routing, acknowledge, and per-target split detection (Sections 8.3 to 8.5).}{SystemVerilog}

\clearpage
\section{Timing Constraints}
\srcentry{Sytembus.sdc}{Sytembus.sdc}{Clock definition and the reset false path quoted in Section 10.}{tcl}

\clearpage
\section{Simulation and Verification}
\srcentry{final_tb.sv}{final\_tb.sv}{System testbench of Sections 12.2 to 12.8: reset, single-master transfers, arbitration, split, and decode error.}{SystemVerilog}
\srcentry{arbiter_tb.sv}{arbiter\_tb.sv}{Unit-level arbiter checks (Section 12.1).}{SystemVerilog}
\srcentry{addr_decoder_tb.sv}{addr\_decoder\_tb.sv}{Unit-level address decoder checks (Section 12.1).}{SystemVerilog}
\srcentry{master2_slave3_tb.sv}{master2\_slave3\_tb.sv}{Two-master and split-capable slave integration test (Sections 12.5 and 12.6).}{SystemVerilog}
\srcentry{remote_link_tb.sv}{remote\_link\_tb.sv}{Bridge round trip over the UART link (Section 4).}{SystemVerilog}
\srcentry{frames_tb.sv}{frames\_tb.sv}{Wire-format checks against the frames of Tables 14 to 16.}{SystemVerilog}
\srcentry{figures_tb.sv}{figures\_tb.sv}{Generates the waveform time windows recorded in figures.txt (Section 12.2).}{SystemVerilog}

\clearpage
\section{Build and Hardware Test Scripts}
\srcentry{scr.tcl}{scr.tcl}{ModelSim driver: compiles the design and runs each testbench (Section 12).}{tcl}
\srcentry{issp_bus_test.tcl}{issp\_bus\_test.tcl}{Scripted hardware verification over JTAG (Section 11).}{tcl}
\srcentry{issp_console.tcl}{issp\_console.tcl}{Interactive memory console over In-System Sources and Probes.}{tcl}
